\section{Four-Valued Probabilistic Epistemic Semantics}
\label{sec:kernel}

The formal development starts with a lightweight, finite-accessibility interface
\[
  M=(W,R,\mu,v,c),
\]
where \(W\) is the world type, \(R_i(w)\) is a finite accessibility list, \(\mu_{i,w}\) is a rational-valued function on event lists, \(v\) assigns FDE values to atoms, and \(c_i\) satisfies
\[
  \frac12<c_i\le 1.
\]
The interface requires only \(\mu_{i,w}(R_i(w))=1\) and
\(\mu_{i,w}(\varnothing)=0\), besides the threshold bounds. It does \emph{not}
require nonnegativity, additivity, extensionality, or duplicate-free lists.
Consequently a general \texttt{Model} is not yet a certified probability model;
``support mass'' below means its rational set-function value unless integrity
is explicitly assumed. The stronger \texttt{StrongProbabilityModel} contract
in Section~\ref{sec:model-paths} supplies ordinary finite-measure laws on
accessible, duplicate-free events. Both signed masses then lie in \([0,1]\).

The stored \texttt{worlds} list need not exhaust the type \(W\), and the interface
does not enforce accessibility closure into that list. The displayed finite
witnesses do use exhaustive finite types and closed ranges. Generic theorems
quantify over the literal Lean interface, not over an unstated finite-universe
axiom. No reflexivity, transitivity, or Euclidean condition is generally assumed.
Normalization does exclude an empty accessible range: otherwise its mass would
be both zero and one (\texttt{PaperReview.accessibility\_nonempty}).

\subsection{Bilateral four-valued information}

A semantic value is a pair of Boolean coordinates
\[
  x=(x^+,x^-),
\]
with the four canonical values
\[
  \Tval=(1,0),\qquad
  \Fval=(0,1),\qquad
  \Bval=(1,1),\qquad
  \Nval=(0,0).
\]
The positive coordinate records support for the formula and the negative coordinate records support against it. Thus \(\Bval\) is a glut and \(\Nval\) is a gap. Negation exchanges the two coordinates, while conjunction takes conjunction on positive support and disjunction on negative support. Independence here means that the coordinates are not logical complements; it does not assert stochastic independence of their support events. In an integrity-certified model the two events share one measure and may overlap.

This bilateral representation is crucial for the dynamic analysis. The two coordinates can cross a probabilistic threshold at different times, producing a temporal structure that would be invisible in a single Boolean belief status.

\subsection{Threshold belief}

For a modal formula \(\varphi\), define its positive and negative accessible extensions at \((i,w)\) by
\[
  E^+_{i,w}(\varphi)
  =\{u\in R_i(w):\llbracket\varphi\rrbracket_u^+=1\},
\]
\[
  E^-_{i,w}(\varphi)
  =\{u\in R_i(w):\llbracket\varphi\rrbracket_u^-=1\}.
\]
The corresponding support masses are
\[
  \Ppos_{i,w}(\varphi)=\mu_{i,w}(E^+_{i,w}(\varphi)),
  \qquad
  \Pneg_{i,w}(\varphi)=\mu_{i,w}(E^-_{i,w}(\varphi)).
\]
Probabilistic belief thresholds the two masses independently:
\[
  \Bel_i\varphi
  =
  \left(
  [\Ppos_{i,w}(\varphi)\ge c_i],
  [\Pneg_{i,w}(\varphi)\ge c_i]
  \right).
\]
Accordingly, \(\Bel_i\varphi=\Bval\) is possible when both signed support masses cross the threshold, while \(\Bel_i\varphi=\Nval\) occurs when neither does.

\subsection{Evidence-stable knowledge}

The modal object language contains primitive probabilistic belief \(\Bel\), primitive evidence-stable knowledge \(\Know\), and primitive raw accessibility possibility. The latter is kept distinct from the internal expression \(\neg\Know\neg\varphi\), because the two need not agree in a four-valued setting.

For the present paper, the central modal predicate is full-value stability. Let
\[
  \Stable_i^M(w,\varphi)
\]
mean pairwise homogeneity on the accessible range:
\[
 H_i(w,\varphi)\iff\forall u,z\in R_i(w),\quad
 \llbracket\varphi\rrbracket_u=\llbracket\varphi\rrbracket_z.
\]
Here \(H_i=\Stable_i^M\). This does not compare the range with the current world
unless that world is accessible. The primitive definition, writing
\(x_u=\llbracket\varphi\rrbracket_u\), is
\[
 (\Know_i\varphi)^+=H_i\land\forall u\in R_i(w)\;x_u^+,
 \qquad
 (\Know_i\varphi)^-=\neg H_i\lor\exists u\in R_i(w)\;x_u^-.
\]
These are classical meta-level Boolean operations on evidence coordinates.
In particular, primitive \(\Know\) does not directly inspect \(\mu\) or \(c\).
Heterogeneity is stipulated to count as negative evidence for this operator.
Without frame assumptions no factivity principle at the current world is
inferred; even on reflexive frames, a glut is not classical truth only.

The resulting relationship between knowledge and belief is exact.

\begin{theorem}[Knowledge as stability-filtered belief]
For every finite 4-PEL model, agent, world, and modal formula,
\[
  \Know_i\varphi
  =
  \begin{cases}
    \Bel_i\varphi,&\Stable_i(\varphi),\\
    \Fval,&\neg\Stable_i(\varphi).
  \end{cases}
\]
\statusverified
\end{theorem}

\begin{proof}[Proof sketch]
The accessible range is nonempty. If homogeneous, each signed extension is
either the entire range or empty. Normalization and \(0<c_i\le1\) make its
threshold bit equal the corresponding common evidence bit. The definition of
\(\Know\) returns those same bits. If heterogeneous, its positive bit is false
and its negative bit true. No additivity assumption is needed.
\end{proof}

Hence an unstable accessible profile is not merely low-confidence evidence. It is a structural failure mode that absorbs every probabilistic belief value into strict \(\Fval\) at the knowledge layer. Conversely, on a homogeneous profile the stability filter is transparent and \(\Know_i\varphi\) preserves the complete four-valued belief state.

For a fixed interpretation \(u\mapsto x_u\) and fixed accessible range, changing
only probability cannot change \(\Know\). This is verified separately by
\texttt{PaperReview.knowledge\_ignores\_probability\_for\_fixed\_values}.
Probability affects \(\Know\varphi\) only through probability-sensitive
subformulas of \(\varphi\); the two-case factorization is not a claim of two
independent numerical inputs to primitive knowledge.
